\documentclass{article}
\usepackage[utf8]{inputenc}
\usepackage{hyperref}
\usepackage{listings}
\usepackage{xcolor}
\usepackage{enumitem}

\lstset{
    basicstyle=\ttfamily\small,
    breaklines=true,
    frame=single,
    backgroundcolor=\color{gray!5},
}

\title{Find large files recursively in a directory tree}
\author{marcos}
\date{2026-04-21}

\begin{document}

\section*{Find large files recursively in a directory tree}
\textbf{Author:} marcos \hfill \textbf{Date:} 2026-04-21

\noindent Recursively finds all regular files larger than 100 MB in the directory tree, displays them with human-readable sizes, and sorts by size in descending order.

\section*{Language}
bash

\section*{Category}
filesystem

\section*{Command}
\begin{lstlisting}
find . -type f -size +100M -exec ls -lh {} + | awk '{print $5, $9}' | sort -hr
\end{lstlisting}

\section*{Explanation}
The command uses 'find' to recursively traverse the directory tree, filter for regular files (-type f), and select those exceeding 100 MB (-size +100M). The '-exec ls -lh \{\} +' displays detailed information for each file in a batch operation, which is more efficient than using '\textbackslash\{\};'. The output is piped to 'awk' to extract the file size and path, then 'sort -hr' sorts the results by human-readable sizes in descending order, showing the largest files first.

\section*{Tags}
find, filesystem, disk-usage, large-files, recursive, storage

\section*{Dependencies}
findutils, coreutils


\section*{Arguments}
\begin{enumerate}
  \item \textbf{PATH} (Optional): Root directory to start the recursive search (default: current directory) \\ Default: .
  \item \textbf{SIZE} (Optional): Minimum file size threshold (default: 100M, supports K, M, G suffixes) \\ Default: 100M
\end{enumerate}

\section*{Examples}
\begin{enumerate}
  \item \texttt{find . -type f -size +100M -exec ls -lh \{\} + | awk '\{print \$5, \$9\}' | sort -hr} - Find files larger than 100 MB in current directory and subdirectories
  \item \texttt{find /var -type f -size +500M -exec ls -lh \{\} + | awk '\{print \$5, \$9\}' | sort -hr} - Find large files in /var directory
  \item \texttt{find . -type f -size +1G -exec ls -lh \{\} + | awk '\{print \$5, \$9\}' | sort -hr} - Find files larger than 1 GB
  \item \texttt{find . -type f -size +50M -exec ls -lh \{\} + | awk '\{print \$5, \$9\}' | sort -hr | head -20} - Show top 20 largest files over 50 MB
\end{enumerate}

\section*{Output}
\begin{lstlisting}
1.5G ./videos/archive.tar.gz
987M ./backups/database.sql.bz2
756M ./datasets/training-data.bin
543M ./cache/node_modules.tar
387M ./downloads/ubuntu-22.04-iso
256M ./logs/app-2026.log
198M ./tmp/large-temp-file
\end{lstlisting}

\section*{Notes}
\begin{itemize}
  \item The -size option supports various units: c (bytes), k (kilobytes), M (megabytes), G (gigabytes), and b (512-byte blocks)
  \item By default, -size rounds up (e.g., +100M matches files larger than 100 MiB)
  \item Use -size -100M to find files smaller than 100 MB
  \item The command may take considerable time on large directory trees
  \item Using '\{\}  +' instead of '\{\}  \textbackslash\{\};' is more efficient as it batches the command executions
  \item Results may be incomplete in directories where you lack read permissions
\end{itemize}

\section*{Warnings}
\begin{itemize}
  \item On network filesystems (NFS, SMB), this command may be significantly slower
  \item Does not follow symbolic links by default; add -L flag if needed
  \item Very large directory trees may cause the command to run for extended periods
  \item Some filesystems may not report accurate sizes for sparse files
\end{itemize}

\section*{See Also}
\begin{itemize}
  \item find-largest-storage-users
  \item disk-space-sort-largest-directories
\end{itemize}

\section*{Status}
reviewed

\section*{Safety}
safe

\section*{Shell}
bash

\section*{Platforms}
linux, gnu-linux, freebsd, openbsd, netbsd

\section*{Created At}
2026-04-21

\section*{Updated At}
2026-04-21

\end{document}
